\subsection{域}

定义 1. 环 K 称为域，如果它不仅仅由 0 组成，并且 K 的每个非零元素都可逆。

域K的非零元素集合在乘法下构成一个群，它正是环K的乘法群 \(K^{*}\) （ \(\S\) 8，第1号）。一个域的反环仍是域。若一个域的乘法是交换的，则称该域为交换域；这样的域与其反环等同。非交换域有时称为除环（体）。

例子。*(1) 我们将在第4号中定义有理数域；在《一般拓扑学》中将定义实数域（《一般拓扑学》，IV，§1，第3号）、复数域（《一般拓扑学》，VIII，§1，第1号）以及四元数域（《一般拓扑学》，VIII，§1，第4号）。除四元数域外，这些域都是交换的。* (2) 环Z/2Z显然是一个域。

设 K 为一个域。K 的每个本身是域的子环 L 称为 K 的子域，此时 K 称为 L 的扩域；这等价于说 L 是 K 的子环，并且对 L 的每个非零元素 x 有 \(x^{-1} \in L\) 。若 \((\mathbf{L}_{i})_{i \in I}\) 是K的一族子域，则 \(\bigcap_{i \in I} L_{i}\) 是 K 的一个子域；因此，对于 K 的每个子集 X，存在包含 X 的 K 的最小子域；称它为由 X 生成的。

命题1. 设 K 是一个域。对于 K 的每个子集 X，X 的中心化子（§8，第5号，例3）X' 是 K 的一个子域。

我们知道（同上引文） \(\mathbf{X}'\) 是 \(\mathbf{K}\) 的子环。另一方面，若 \(x \neq 0\) 与 \(z \in \mathbf{X}\) 可交换，则 \(x^{-1}\) 亦然（I，§2，第3号，命题5），从而 \(\mathbf{X}'\) 包含 \(\mathbf{X}'\) 中每个非零元素的逆。

推论。域K的中心是K的一个（交换）子域。

定理1. 设A是一个环。以下条件等价：

\begin{enumerate}
\def\labelenumi{(\alph{enumi})}
\item
  A 是一个域；
\item
  A 不约化为 0，且 A 仅有的左理想是 0 和 A。
\end{enumerate}

假设A是一个域。那么A不约化为0。设 a 是 A 的一个异于 0 的左理想。存在一个非零的 a 属于 a。对于所有 \(x \in A\) , \(x = (xa^{-1})a \in a\) ；因此 a = A。

假设 A 满足条件 (b)。设 \(x \neq 0\) 属于 A。我们需要证明 x 是可逆的。左理想 Ax 包含 x，因此不为零，从而 Ax = A。于是存在 \(x' \in A\) 使得 \(x'x = 1\) 。则 \(x' \neq 0\) ，因为 \(1 \neq 0\) 。将上述结果应用于 \(x'\) ，可以看出存在 \(x'' \in A\) 使得 \(x''x' = 1\) 。则 \(x'' = x'' \cdot 1 = x''x'x = 1 \cdot x = x\) ，因此 \(xx' = 1\) 。因此 \(x'\) 是 x 的逆。

注。在定理1中，左理想可以换成右理想。在第八章，§5，第2号中，我们将研究没有异于0和A的双边理想的非零环A；这样的环（称为拟单环）不一定是域 *（例如，以有理数为系数的2阶方阵组成的环 \(\mathbf{M}_2(\mathbf{Q})\) 是拟单的，但不是域）*。

推论1. 设A是一个环，且a是A的一个双边理想。要使环A/a成为域，当且仅当a是A的一个极大左理想。

A/a的左理想具有b/a的形式，其中b是A的包含a的左理想（§8，第8号，命题5的推论）。说A/a ≠ 0意味着a ≠ A。在此假设下，A/a是域当且仅当A的包含a的左理想仅有a和A（定理1），由此得出推论。

推论2. 设A是一个非零交换环。则存在一个从A到某个交换域上的同态。

由 Krull 定理（§ 8，第 6 号，定理 1），在 A 中存在一个极大理想 \(a\) 。则 A/a 是一个域（推论 1）。

推论 3. 设 a 为整数 \(\geqslant 0\) 。环 \(\mathbf{Z} / a\mathbf{Z}\) 为域当且仅当 a 为素数。

这由推论1及 §8 第11号得出。

对于素数p，域Z/pZ记作 \(F_{p}\) .

定理2. 设K是一个域，A是一个非零环。若f是K到A中的一个同态，则A的子环 \(f(\mathbf{K})\) 是一个域，且f定义了K到 \(f(\mathbf{K})\) 上的一个同构。

设 \(\mathfrak{a}\) 为 \(f\) 的核。则 \(1 \notin \mathfrak{a}\) ，因为 \(f(1) = 1 \neq 0\) 在 \(\mathbf{A}\) 中，并且由于 \(\mathfrak{a}\) 是 \(\mathbf{K}\) 的左理想，由定理1有 \(\mathfrak{a} = \{0\}\) 。因此 \(f\) 是单射，从而是 \(\mathbf{K}\) 到子环 \(f(\mathbf{K})\) 上的同构，该子环位于 \(\mathbf{A}\) 中；因此后一环是一个域。

